\section{The Part II convexity construction}
\label{sec:annex-p17-astrom-incomplete-information-ii}

\begin{annexitem}{annex:p17.astrom.ii.mass.bridge}{thm:p17-astrom-ii-mass-bridge}{Positive observation maps and the exact probability interface}
\begin{proof}
All entries of \(A_{j,t}(u)\) are nonnegative.
For \(w\in S\), its image is therefore in \(S\), and
\[
 \begin{split}
 \sum_j\|A_{j,t}(u)w\|_1
 &=\sum_{j,i,s}q_{ij,t}(u)p_{si,t}(u)w_s\\
 &=\sum_s w_s\sum_i p_{si,t}(u)\sum_jq_{ij,t}(u)
 =\sum_s w_s=\|w\|_1 .
 \end{split}
\]
Both stochastic row-sum identities are used.
This is the extension of the 1965 unnormalized
posterior numerator to control-dependent observations.
At \(b\in\Delta_n\), its norm is the probability of
the next observation \(j\), and its normalized image
is the next posterior on every positive-probability
branch. The proof of \cref{thm:p17-astrom-filter}
applies with \(q_{ij}\) replaced by the specified
\(q_{ij,t}(u)\), conditional on the known chosen action.
The sum above shows that the observation probabilities
still sum to one; zero image means a zero-probability
branch. Borel entries give all the same kernel
measurability and default conventions.

At \(N\), \(|\langle g_N(u),w\rangle|
\leq G_N\|w\|_1\), uniformly in \(u\).
Suppose the next value has absolute value at most
\(C_{t+1}\) on \(\Delta_n\).
The continuation at \(w\) then has absolute value
at most
\(C_{t+1}\sum_j\|A_{j,t}(u)w\|_1
 =C_{t+1}\|w\|_1\).
Adding the immediate term and taking either extremum
gives \(|V_t^\pm(w)|\leq C_t\|w\|_1\).
The extrema are finite because \(U\) is nonempty
and this two-sided bound applies to every candidate.

For \(c>0\), linearity gives
\(A_{j,t}(u)(cw)=cA_{j,t}(u)w\).
Its norm scales by \(c\), and its normalized nonzero
image does not change. Every term in the objective
therefore scales by \(c\), as does either extremum.
At \(c=0\), all images and the immediate reward
vanish, giving the declared value zero.
This proves positive homogeneity. Applying it to
each nonzero \(A_{j,t}(u)w\) changes the weighted
normalized term into \(V_{t+1}^\pm(A_{j,t}(u)w)\);
when the image is zero both versions contribute
zero. This proves the unnormalized recursion.
Taking \(c=\|w\|_1\) and normalized input
\(w/\|w\|_1\) proves the cone/simplex identity.

\paragraph{Equality with the control value without assuming attainment.}
At a fixed starting stage, consider the same class
of Borel observation/action-history policies for
every possible starting hidden state.
For a fixed policy \(\pi\), let \(a_i^\pi\) be its
expected remaining reward conditional on starting
in hidden state \(i\), before any additional
observation. The finite kernel product law gives
this number, with \(|a_i^\pi|\leq C_t\).
Starting from belief \(b\) gives value
\(\sum_i b_i a_i^\pi\), by linearity of the initial
mixture and the kernel integrals.
Although a controller may select a policy after
learning its prior, the available class of all
such history kernels is the same at every prior.

Let \(W_t^+(b)\) be the supremum of these policy
values, and \(W_t^-(b)\) the infimum.
These are bounded functions. They are Borel because
they are Lipschitz:
for every policy
\[
 \left|\sum_i(b_i-b'_i)a_i^\pi\right|
                              \leq C_t\|b-b'\|_1,
\]
and taking suprema or infima preserves this same
bound. This argument does not assume the policy
family or its coefficient vectors to be finite.

At the terminal date, randomizing a control cannot
exceed the supremum, or fall below the infimum,
of its immediate expected rewards. Conversely a
constant control approximating that extremum is
an admissible policy. Hence \(W_N^\pm=V_N^\pm\).
Suppose \(W_{t+1}^\pm=V_{t+1}^\pm\).
For an arbitrary first action and observation,
the conditional future history policy is admissible
from the resulting posterior. Its value is bounded
by \(W_{t+1}^+\) in the maximum case and by
\(W_{t+1}^-\) in the minimum case.
The Borel property just established permits
integration of these bounds against a randomized
first-action kernel. Averaging cannot improve
on the corresponding extremum over first controls.
This proves the upper bound for the maximum
problem and the lower bound for the minimum problem.

For the reverse maximum inequality at a fixed
belief \(b\), choose a first control within
\(\varepsilon/2\) of the supremum in the recursion.
For each of the finitely many positive-probability
observation branches choose a continuation policy
within \(\varepsilon/2\) of its value.
Use arbitrary defaults on the other branches.
Pasting this finite family by the first observation
is a Borel history policy. The weighted sum of its
continuation losses is at most \(\varepsilon/2\),
so its total value is within \(\varepsilon\)
of the recursive supremum.
For the infimum choose controls and continuation
policies within the same tolerances from above.
Let \(\varepsilon\downarrow0\).
This proves \(W_t^\pm=V_t^\pm\) by induction,
without asserting that a supremum is attained.
The approximating policies can be chosen deterministic:
the terminal approximation is a constant control, and
the induction pastes finitely many deterministic
continuations after a deterministic first control.
Thus the original deterministic history class and
the larger randomized class have the same extremal
values, whether or not those values are attained.
When \(Q_t\) is independent of the control,
all numerator, likelihood, posterior and
Bellman formulas specialize literally to the
1965 interface.
\end{proof}
\end{annexitem}

\begin{annexitem}{annex:p17.astrom.ii.maxima}{thm:p17-astrom-ii-maxima}{The four scalar cases and an arbitrary family}
\begin{proof}
Here are the four cases in the scalar inequality,
with equality boundaries included.
If \(a\geq c\) and \(b\geq d\), the right side is
\(a+b\), which is at least both left candidates.
If \(a\geq c\) and \(b\leq d\), the right side is
\(a+d\), at least \(a+b\) and \(c+d\).
If \(a\leq c\) and \(b\geq d\), it is \(c+b\),
again at least both. If \(a\leq c\) and \(b\leq d\),
it is \(c+d\), at least both.
The cases cover all inputs; their possible overlap
on ties gives the same inequality.

For convex \(f_1,f_2\), write
\(f=\max(f_1,f_2)\) and \(z=\theta x+(1-\theta)y\),
where \(0\leq\theta\leq1\).
Convexity followed by the scalar inequality gives
\[
 \begin{split}
 f(z)&\leq
 \max\{\theta f_1(x)+(1-\theta)f_1(y),
                \theta f_2(x)+(1-\theta)f_2(y)\}\\
 &\leq\theta\max(f_1(x),f_2(x))
             +(1-\theta)\max(f_1(y),f_2(y))\\
 &=\theta f(x)+(1-\theta)f(y).
 \end{split}
\]
The second maximum in the middle line is evaluated
at \(y\). Nonnegative coefficients allow them to
be taken outside the maxima, including zero
coefficients at the mixture endpoints.

For any nonempty family \((f_\alpha)\) with finite
pointwise supremum \(F\), each member satisfies
\[
 f_\alpha(z)\leq
       \theta f_\alpha(x)+(1-\theta)f_\alpha(y)
                \leq\theta F(x)+(1-\theta)F(y).
\]
Take the supremum on the left.
No maximizing index or enumeration of the family
is needed. For finite sums, add the convexity
inequalities term by term.
For concave functions replace every function
by its negative; a negative infimum is the
supremum of the negatives. This proves the
corresponding infimum and sum statements.
\end{proof}
\end{annexitem}

\begin{annexitem}{annex:p17.astrom.ii.perspective}{thm:p17-astrom-ii-perspective}{Every norm and zero-image case in Lemma 2}
\begin{proof}
For \(x,y\in S\), positivity of \(A\) gives
\(Ax,Ay\in S\). On this cone the chosen norm
is additive and positively homogeneous:
\[
 \|\theta Ax+(1-\theta)Ay\|_1
          =\theta\|Ax\|_1+(1-\theta)\|Ay\|_1.
\]
This is the sum of the coordinate equalities,
which hold because every coordinate is nonnegative.
It is not a norm identity for arbitrary signed
vectors or arbitrary norms.

If \(\theta=0\) or \(1\), the required convexity
inequality is equality. Suppose \(0<\theta<1\).
Put \(a=\|Ax\|_1\), \(b=\|Ay\|_1\).
If \(a=b=0\), all three images are zero and all
three perspective values are zero.
If \(a=0<b\), then \(Ax=0\) and
\[
 A(\theta x+(1-\theta)y)=(1-\theta)Ay.
\]
Its perspective value is exactly
\((1-\theta)({\cal P}_Af)(y)\), which is the
required convex combination with
\(({\cal P}_Af)(x)=0\).
The case \(b=0<a\) is the symmetric equality case.

For \(a,b>0\), let \(c=\theta a+(1-\theta)b>0\).
Then
\[
 \frac{A(\theta x+(1-\theta)y)}c
       =\frac{\theta a}{c}\frac{Ax}{a}
             +\frac{(1-\theta)b}{c}\frac{Ay}{b}.
\]
The two coefficients are nonnegative and sum to
one, and both normalized images lie in \(\Delta_n\).
Convexity of \(f\), multiplied by \(c\), therefore
gives
\[
 ({\cal P}_Af)(\theta x+(1-\theta)y)
      \leq\theta a f(Ax/a)+(1-\theta)b f(Ay/b)
      =\theta({\cal P}_Af)(x)+(1-\theta)({\cal P}_Af)(y).
\]
This proves convexity on the full closed cone.
For any \(r>0\), linearity of \(A\) and norm
homogeneity give
\(({\cal P}_Af)(rx)=r({\cal P}_Af)(x)\);
when \(r=0\) both sides are zero by definition.
For concave \(f\), apply the proved result to
\(-f\) and use
\({\cal P}_A(-f)=-{\cal P}_Af\).
All zero cases remain valid, so the concave
conclusion has the same full domain.
\end{proof}
\end{annexitem}

\begin{annexitem}{annex:p17.astrom.ii.convexity}{thm:p17-astrom-ii-convexity}{The full source induction for maximum reward}
\begin{proof}
For each fixed terminal control \(u\),
\(w\mapsto\langle g_N(u),w\rangle\) is linear
and therefore convex on \(S\).
Its finite-valued supremum over the common control
set is convex by the arbitrary-family lemma.
This proves the terminal case without treating
an uncountable control set as a finite list.

Suppose \(V_{t+1}^+\) is convex on \(S\).
For a fixed control \(u\) and observation \(j\),
the map \(A_{j,t}(u)\) is positive and linear.
The perspective lemma therefore proves convexity of
\[
       w\longmapsto
       \begin{cases}
       \|A_{j,t}(u)w\|_1
          V_{t+1}^+(A_{j,t}(u)w/\|A_{j,t}(u)w\|_1),
                                 &A_{j,t}(u)w\ne0,\\
       0,&A_{j,t}(u)w=0.
       \end{cases}
\]
Their finite sum is convex. Adding the linear
immediate reward preserves convexity.
Finally take the supremum over the same \(U\)
at every \(w\); the finite-valued supremum lemma
proves convexity of \(V_t^+\).
Backward induction reaches every \(t\).
The mass bound already proved guarantees finiteness
at every step. If the supremum is attained, it
is the source's maximum with no change in this proof.
Only the dependence on \(w\) is required to be
linear for each fixed control; no linearity
in \(u\) was used.
For the broader algebraic formulation, replace the
mass-bound justification of finiteness by its explicit
finite-value hypothesis. Every other step is unchanged:
positive linear maps, finite sums and an arbitrary
common control supremum are exactly the operations
already proved to preserve convexity.
\end{proof}
\end{annexitem}

\begin{annexitem}{annex:p17.astrom.ii.concavity}{thm:p17-astrom-ii-concavity}{Every step of the minimum-cost dual}
\begin{proof}
At the terminal date,
\[
 \inf_u\langle g_N(u),w\rangle
       =-\sup_u\langle-g_N(u),w\rangle.
\]
Suppose the analogous identity
\(V_{t+1}^-[g]=-V_{t+1}^+[-g]\) holds.
Substitute it in every positive-weight continuation
term of the minimum recursion. Its candidate
objective is the negative of the corresponding
maximum candidate for reward vector \(-g_t\).
The identity \(\inf_u(-a_u)=-\sup_u a_u\)
then proves \(V_t^-[g]=-V_t^+[-g]\).
Zero observation columns have zero contribution
on both sides; the origin and terminal conventions
also agree. Induction proves the identity at
all stages.

The maximum theorem applied to \(-g\) gives
\[
 V_t^+[-g](\theta w+(1-\theta)v)
       \leq\theta V_t^+[-g](w)
                   +(1-\theta)V_t^+[-g](v).
\]
Multiply by \(-1\) and use the proved identity
to obtain the required concavity inequality.
Equivalently, one can apply the concave
perspective lemma, add its finitely many terms,
and take the infimum of the resulting concave
candidates at each step.
The actual Part I switch criterion minimizes
terminal failure probability, whereas the
production criterion maximizes profit. Their
directions are therefore exactly as stated.
\end{proof}
\end{annexitem}

\begin{annexitem}{annex:p17.astrom.ii.continuity}{thm:p17-astrom-ii-continuity}{The full Lipschitz bound and the controlled-observation extension}
\begin{proof}
For every \(w,v\in S\) and fixed \(u,t\),
the coordinate triangle inequality gives
\[
 \begin{split}
 \sum_j\|A_{j,t}(u)w-A_{j,t}(u)v\|_1
 &\leq\sum_{j,i,s}q_{ij,t}(u)p_{si,t}(u)|w_s-v_s|\\
 &=\sum_s|w_s-v_s|=\|w-v\|_1 .
 \end{split}
\]
At \(N\), every linear reward candidate has
Lipschitz constant at most \(G_N\) in this norm.
The supremum and infimum of functions with a
common Lipschitz constant retain that constant:
each \(f_\alpha(w)\leq f_\alpha(v)+L\|w-v\|_1\)
can be bounded by the common supremum at \(v\),
then maximized; reverse \(w,v\).
The infimum case follows by negation.

Use the unnormalized recursion and suppose the
next value is \(C_{t+1}\)-Lipschitz on \(S\).
Each candidate has difference at \(w,v\) bounded
in absolute value by
\[
 G_t\|w-v\|_1+
 C_{t+1}\sum_j\|A_{j,t}(u)w-A_{j,t}(u)v\|_1
                      \leq C_t\|w-v\|_1.
\]
Taking either extremum proves the bound at \(t\).
Induction proves the full claim, including
the origin and every face, without assuming
continuity of \(Q_t\) in the control.

Now impose that additional continuity.
Every coordinate of \(A_{j,t}(u)b\), its mass,
and the immediate reward is continuous in
\((b,u)\in\Delta_n\times U\).
The next value is continuous and bounded on
the simplex by the proved bound.
The weighted term
\(\|A_{j,t}(u)b\|_1
 V_{t+1}^\pm(A_{j,t}(u)b/\|A_{j,t}(u)b\|_1)\)
is continuous where its mass is positive.
At zero mass its absolute value is at most
\(C_{t+1}\|A_{j,t}(u)b\|_1\), tending to zero.
Thus the candidate objective is jointly continuous
on \(\Delta_n\times U\).

The complete compact-minimizer construction
\cref{thm:p17-astrom-continuous-selection} now
applies with precisely its stated hypotheses.
For a maximum apply its minimizing selector to
the negative objective.
Its nested closed-cube construction gives a
Borel optimum at every belief, including all
boundary beliefs. No continuous selector is
presumed. Backward induction gives deterministic
belief policies attaining the value at every
starting stage and belief.
The filter proof uses the known chosen \(u\)
in both \(P_t\) and \(Q_t\); conditioning on
the private action does not change the
current hidden-state distribution.
Therefore its numerator and likelihood are
exactly \(A_{j,t}(u)b\) and its norm.
This supplies the conditional law needed by
the full backward policy argument in
\cref{thm:p17-astrom-finite-optimum}.

For clarity, the other source interfaces continue
for concrete reasons. Implementation computes the
same posterior from the observed history and known
actions. Markov projection disintegrates the action
given that posterior under a fixed initial law,
then uses this common controlled belief transition.
The conditional cost identity is unchanged.
The Bellman residual proof still telescopes these
conditional cost and next-value identities, with
the terminal value zero.
In the perfect-information bound,
\(\sum_j\|A_{j,t}(u)b\|_1F_j(i)
 =\sum_s b_sp_{si,t}(u)\), so the same
statewise lower bound follows after summation.
The no-further-observation bound still follows
by restricting to a fixed control schedule.
These directions are for the minimum-cost criterion:
perfect information is a lower bound and the schedule
value is an upper bound. Negating rewards proves the
reverse directions for maximum reward, with perfect
information above and the schedule value below.
If the rows of \(Q_t(u)\) are identical for each
fixed \(t,u\), all its positive-probability
posteriors equal \(bP_t(u)\), so its optimum
is that schedule value. The common row may
depend on the known control.
If \(Q_t(u)=I\) for all \(t,u\), the positive
posteriors are point masses, giving the
complete-state recursion. Initial comparisons
use the separately specified initial observation.

Finally an optimum \(f_t\) on \(\Delta_n\)
extends to the cone by
\(\widehat f_t(w)=f_t(w/\|w\|_1)\) for \(w\ne0\)
and a fixed control at \(0\).
This is Borel, constant on positive rays,
and optimizes there by positive homogeneity.
The discontinuous-observation counterexample
below proves why convexity and compact \(U\)
alone do not justify the attainment step.
\end{proof}
\end{annexitem}

\begin{annexitem}{annex:p17.astrom.ii.examples}{thm:p17-astrom-ii-examples}{The actual Part I examples under homogeneous extension}
\begin{proof}
The full four-stage switch calculation in
\cref{thm:p17-astrom-example-one} gives, at
normalized belief \((x,1-x)\) and \(q=4/5\),
\[
 V_1^-(x,1-x)=
 \min\{x,1-x,\tfrac9{25}x+\tfrac1{25}(1-x),
                         \tfrac9{25}(1-x)+\tfrac1{25}x\}.
\]
For \(w\ne0\), put \(x=w_1/(w_1+w_2)\)
and multiply by \(w_1+w_2\).
The four terms become exactly the four linear
forms in the stated cone expression. At
\(w=0\) each form is zero, agreeing with
the prescribed extension.
At \((2,1)\) the candidates are
\(2,1,19/25,11/25\), so the minimum is \(11/25\).
At the two stated normalized beliefs, the
smaller affine form is \(9/125\).
At their midpoint \((1/2,1/2)\) the two
noncoordinate forms are \(1/5\), smaller
than \(1/2\). Thus the midpoint value
strictly exceeds the average \(9/125\),
which is the concavity direction.

For production, the fully calculated
\cref{thm:p17-astrom-example-two} gives
\(V_3^+(x,1-x)=\max(27.5x,20x+4.5)\).
The same homogeneous extension changes its
second candidate into
\(24.5w_1+4.5w_2\), while the first becomes
\(27.5w_1\). At \((2,1)\) they are
\(55\) and \(53.5\), so the maximum is \(55\).
At \((1,0)\) the value is \(27.5\), and
at \((0,1)\) it is \(4.5\).
At their midpoint the two candidates are
\(13.75\) and \(14.5\), so the value is
\(14.5\), strictly below their average \(16\).
This is the convexity direction.
The formulas are maxima or minima of
homogeneous linear forms, so their origin
and scaling assertions hold directly as well.
These computations concern exactly the two
finite examples; they do not establish a
finite piecewise-linear representation for
an arbitrary continuous control family.
\end{proof}
\end{annexitem}

\begin{annexitem}{annex:p17.astrom.ii.boundaries}{thm:p17-astrom-ii-boundaries}{Five explicit failures when a premise is removed}
\begin{proof}
\paragraph{A nonlinear cone map.}
Take \(S=\mathbb R_+\), one control and one
observation term, zero current reward, and
terminal value \(V_N(w)=w\).
Replace the required linear map by \(A(w)=\sqrt w\).
It maps \(S\) into \(S\), but the preceding
value is
\(\sqrt w\,V_N(1)=\sqrt w\), with value zero
at \(w=0\).
At \(w=0,1\) the values average to \(1/2\);
at their midpoint the value is \(1/\sqrt2>1/2\).
Thus mere mapping of \(S\) into \(S\)
does not give the theorem.

\paragraph{A different norm.}
On \(S=\mathbb R_+^2\), take \(A\) to be
the identity and the everywhere constant
convex function \(f=-1\).
Replacing \(\|\cdot\|_1\) by the Euclidean
norm makes its weighted normalized value
\(-\|w\|_2\). At the two coordinate unit
vectors this equals \(-1\); at their
midpoint it equals \(-1/\sqrt2>-1\).
The resulting function is not convex.
The particular additive norm on the
nonnegative cone is therefore essential
to the source's weight calculation.

\paragraph{A signed linear map.}
Let \(A(w_1,w_2)=(w_1-w_2,0)\).
It does not map \(S\) into \(S\), so the
source's normalized argument can leave
its domain. Even extending \(f=-1\)
to all of \(\mathbb R^2\) does not repair
convexity: the perspective becomes
\(-|w_1-w_2|\), zero on equal coordinates.
At the two unit vectors its values are
\(-1,-1\), while at their midpoint it
is zero. This violates the convexity
inequality.

\paragraph{Belief-dependent feasible controls.}
For a one-stage normalized two-state problem,
let action zero have reward zero in both
states and action one have reward one
in both states.
Permit only action zero at beliefs with
\(b_1<1/2\), and only action one otherwise.
The maximum value is then zero below
that threshold and one at or above it.
At beliefs \((.49,.51)\) and \((.51,.49)\)
the values are zero and one; at their
midpoint the value is one, exceeding
their average \(1/2\).
There is no common family of candidate
functions available at every belief,
which is the premise used in the supremum lemma.

\paragraph{Discontinuous controlled observations and nonattainment.}
Use two stages, a hidden fair bit which stays
fixed, and compact controls \(U=[0,1]\).
The first reward is zero.
At the second stage, reward is \(u\) in
state one and \(1-u\) in state two.
These rewards and the identity transition
matrix are continuous in \(u\).
Let the first chosen control determine
the accuracy
\[
 h(u)=
 \begin{cases}(1+u)/2,&0\leq u<1,\\1/2,&u=1,\end{cases}
 \qquad
 Q(u)=\begin{pmatrix}h(u)&1-h(u)\\1-h(u)&h(u)\end{pmatrix}.
\]
This is a Borel stochastic observation
matrix, discontinuous at \(u=1\).
Given the fair prior and either observation,
the posterior probability of the indicated
state is \(h(u)\geq1/2\).
At the second stage, maximizing the
linear reward selects the more likely
state, giving expected reward exactly
\(h(u)\) after a first deterministic \(u\).
The supremum over first controls is one,
approached as \(u\uparrow1\), but \(h(u)<1\)
for every \(u\in U\).

Randomization does not attain the supremum.
For any first-action law, the nonnegative
variable \(1-h(U_1)\) is strictly positive
almost surely and therefore has strictly
positive expectation. Indeed the sets
\(\{1-h(U_1)>1/k\}\) cover the probability
space, so some has positive probability.
The expected final reward is consequently
at most \(E h(U_1)<1\).
Thus no policy attains the optimum.
For every positive tolerance a constant
\(u<1\) sufficiently near one is within
that tolerance.
More generally the supremum value at
any normalized prior is one: observations
approaching perfect accuracy and then
choosing the indicated state give error
at most \(1-h(u)\), while rewards never
exceed one. Its cone extension is
\(w_1+w_2\), a linear, hence convex, value.
This separates the source's convexity
conclusion from an unjustified existence
claim under discontinuous sensing.
\end{proof}
\end{annexitem}
